\section{Evaluation}\label{sec:evaluation}

We ask five questions. \textbf{RQ1}: does the model reach the target levels (recall and precision $\ge 95\,\%$, latency
p95 $\le 3$\,s) on public datasets? \textbf{RQ2}: which design choices carry the result? \textbf{RQ3}: does it hold on a
live cluster, blind, and at what alert burden? \textbf{RQ4}: can it tell an attack from a harmless new process?
\textbf{RQ5}: what does it cost, and how fast does an alert reach a responder?

\subsection{Methodology}\label{sec:eval-method}

\paragraph{Pre-registration.} Every main analysis and every live run has a JSON protocol committed to version control
before the data it concerns is read: expected levels, metric definitions, selection rules and analysis code. Later
changes are timestamped amendments. We separate \emph{development} results (data seen while designing) from
\emph{blind} results (the official split of LID-DS~2021 and live runs); conclusions rest on blind results only.

\paragraph{Datasets.} LID-DS~2021~\cite{grimmer2022} has 15 scenarios of real vulnerabilities. Each scenario's model
learns its training and validation recordings. The test recordings are split by a registered seed into a development
part (30\,\%: 544 attack and 3{,}403 normal recordings) and an official part (70\,\%: 1{,}271 attack and 7{,}938
normal). An attack recording is detected when an alarm window ends after the first exploit timestamp; latency is
measured from that timestamp; a normal recording with any alarm is a false alarm. CB-DS~\cite{elkhairi2022} has 1{,}700
normal recordings of an e-shop container and nine container-breakout scenarios of 99 recordings each; it has no exploit
timestamp, so any alarm in an attack recording counts and latency is not measured. LID-DS~2019 is reserved for one
final run (\cref{sec:eval-2019}).

\paragraph{Cluster campaign.} Each trial targets one replica with consecutive 15\,s intervals on the same container: an
idle control; a harmless control binary started through the same \texttt{kubectl exec} path as the attack, which only
loops over \texttt{getpid} and \texttt{clock\_nanosleep}; from the second run on, a real operator command
(\texttt{ls}, \texttt{cat}, \texttt{ps}, \texttt{env}, \texttt{id}, \texttt{sh}, \texttt{python}, \texttt{date},
\texttt{uname}, \texttt{df} in turn); and one of five self-written attack fixtures repeated at a registered rate:
a child calling \texttt{ptrace(PTRACE\_TRACEME)}; \texttt{setns} on file descriptor $-1$; \texttt{execveat} of a
missing file; flipping an anonymous page between read and read-write with \texttt{mprotect}; an invalid
\texttt{seccomp} operation. Every action affects only the calling process or fails before changing host state.
21 workloads $\times$ 5 fixtures $\times$ 5 rates (1, 3, 6, 12, 24 actions/s) give 525 trials per run; the
\emph{primary set} is the 315 trials at 6--24/s, rates 1 and 3 model a slow attacker. Trial order is shuffled with a
fixed seed and attacks start after a 90-minute selection prefix.

\paragraph{Metrics.} Recall: an attack interval with an alarm within its 15\,s. Control positives: idle intervals with
an alarm. Precision 1:1: TP/(TP+control positives), since each trial has one attack and one idle interval; harmless
binaries and operator commands are reported apart (\cref{sec:eval-benign}). Latency: from the first kernel event of
the attack to the end of the first alarm window, and separately to Alertmanager. Burden: alarm episodes (alarms of one
container less than 15\,s apart) per replica-hour of normal traffic outside the selection prefix and the campaign.
Rates carry Wilson 95\,\% intervals~\cite{wilson1927}; burden a workload-clustered bootstrap~\cite{efron1993}.

\subsection{RQ1: public datasets}\label{sec:eval-datasets}

The protocol registered three levels, each equal to the best result reached before on the same data. All three are met.
On development the model detects 530/544 recordings (97.43\,\%) with 6/3{,}403 false-alarm recordings and p95
1.52\,s; on CB-DS 891/891 with 3/1{,}410. \Cref{tab:lidds} gives the blind official result.

\begin{table}[t]
\centering\small
\caption{LID-DS 2021 official split (1{,}271 attack, 7{,}938 normal recordings); 95\,\% Wilson intervals.}
\label{tab:lidds}
\begin{tabular}{@{}lll@{}}
\toprule
Metric & Target & Result \\
\midrule
Recall & $\ge 95\,\%$ & 97.40\,\% (1{,}238/1{,}271) [96.38, 98.15] \\
False-alarm recordings & reported & 0.33\,\% (26/7{,}938) [0.22, 0.48] \\
Precision at base rate & $\ge 95\,\%$ & 97.9\,\% \\
Latency p50 / p95 & p95 $\le 3$\,s & 0.39\,s / 1.47\,s \\
Detected within 3\,s & reported & 1{,}234/1{,}271 \\
\bottomrule
\end{tabular}
\end{table}

Fourteen of 15 scenarios reach 100\,\% recall; every miss belongs to Juice-Shop (52/85). Of its 33 misses, 28 leave no
trace in any signal we tried (new path, new error code, size, unseen sequence), which caps those signals at
1{,}243/1{,}271. The missed requests are logins with a modified SQL string to the same endpoint as normal logins: no
new process, path, error or size, and the traces keep only the first ${\sim}80$ bytes of buffers
(\cref{sec:discussion}). False alarms concentrate on one scenario (8/532); all others stay below 0.6\,\%. Precision
depends strongly on the base rate: with the same recall and false-alarm rate it would be 75.0\,\% at 1\,\% attacks and
22.9\,\% at 0.1\,\%~\cite{arp2022}. The only published result on all 15 scenarios is the dataset authors' STIDE, with a
mean per-scenario detection rate of 0.59~\cite{grimmer2022}; computed the same way ours is 0.97 (their false alarms are
counted per alarm cluster, so only detection rates compare). The official split is not entirely unseen: the alphabet
was designed after looking at other signals on the same recordings.

\paragraph{Why tokens.} An earlier version of the system scored per-window system-call count vectors with
ExtraTrees~\cite{geurts2006} and a value-range branch. On all LID-DS~2021 test recordings it detected 1{,}154/1{,}815
(63.6\,\%, p95 4.13\,s); the misses concentrate on exploits that run inside the service process (e.g.\ 113 of 120
Heartbleed and 114 of 118 ZipSlip recordings), where the set and frequency of calls hardly change. On the official split
with the same threshold rule it detects 662/1{,}271 (p95 13.15\,s); next to the token model it adds no detection and
raises false-alarm recordings from 26 to 36.

\subsection{RQ2: which design choices matter}\label{sec:eval-design}

All numbers in this subsection are development and CB-DS; the official split was not read. A cache path reproduces the
stored scores on all 16 scenarios (largest difference $2.8\times10^{-14}$) before any variant is read.

\begin{table}[t]
\centering\small
\caption{Token alphabets on LID-DS 2021 development (544 attack recordings; pre-cross-fit comparison at $\alpha=0.002$).}
\label{tab:alphabets}
\begin{tabular}{@{}lrrr@{}}
\toprule
Alphabet & Detected & FA rate & p95 \\
\midrule
\textbf{select1f} (used) & \textbf{530} & 0.56\,\% & 1.45\,s \\
select1f in probe emission order & 530 & 0.56\,\% & 1.49\,s \\
errno of every call, halved sizes & 527 & 0.59\,\% & 1.52\,s \\
no error codes & 528 & 0.50\,\% & 1.53\,s \\
no open mode and directory & 523 & 0.56\,\% & 1.51\,s \\
requested size at entry & 502 & 0.38\,\% & 2.29\,s \\
no sizes & 490 & 0.56\,\% & 2.58\,s \\
system-call name only & 487 & 0.56\,\% & 2.59\,s \\
\bottomrule
\end{tabular}
\end{table}

\paragraph{Information in a token.} \Cref{tab:alphabets}: read and write sizes carry the most signal, and it must be
the \emph{returned} size, not the requested one, which is why the probe hooks the exit of a call. Names alone lose 40
recordings against the full alphabet. Emitting unhooked calls at entry, as the probe does, loses nothing.

\paragraph{Each of the 39 hooked calls.} Reducing one hooked call at a time to its name and rescoring the whole
pipeline changes the result for only 7 of the 31 calls present in the data: \texttt{write} ($-40$ detections),
\texttt{open} ($-4$), \texttt{read} ($-2$, +1 false alarm), \texttt{pread} ($-1$), and \texttt{recv}/\texttt{recvfrom},
\texttt{send}/\texttt{sendto} and \texttt{stat} only in false alarms. A registered backward elimination (36 scorings)
leaves five calls (\texttt{open}, \texttt{read}, \texttt{recv} with \texttt{recvfrom}, \texttt{stat}, \texttt{write})
that give exactly the result of all 39 in every scenario. Under a declared reuse of the official split, after the full
set reproduced its official result, the five-call set detects the same 1{,}238/1{,}271 with 17 instead of 26
false-alarm recordings (p95 1.46\,s). Adding the details of 13 other calls never increases detections. The 39-call set
is sufficient but wider than this dataset needs; we keep it because narrowing the probe changes what the cluster sees.

\paragraph{Parameters.} Varying one parameter at a time around the chosen values leaves the result flat: gram length 3,
a top-3 to top-5 mean, 3 to 10 folds and $\alpha'$ from $5\times10^{-4}$ to $3\times10^{-3}$ all keep 528--530
detections. It breaks only at the extremes: the top-1 surprise detects 123/544 (one strange event does not separate
attacks from normal), a top-20 mean 492, grams of length 2 or 6 detect 458 and 493.

\paragraph{Is the Gumbel threshold right?} Not as a distribution: an Anderson--Darling test rejects a Gumbel fit of the
held-out maxima in 15/15 scenarios after Holm correction, and the estimated GEV shape ranges from $-0.28$ to $+0.77$.
Nor as a level: at $\alpha'=0.01$ we expect 34 normal recordings above threshold and see 51; at $1.8\times10^{-4}$ we
expect 0.6 and see 6. As a \emph{rule}, however, it is the best normal-only rule we tried (\cref{tab:thresholds}).

\begin{table}[t]
\centering\small
\caption{Most detections of each normal-only threshold rule within a false-alarm budget (normal recordings), development.}
\label{tab:thresholds}
\begin{tabular}{@{}lrrr@{}}
\toprule
Rule & $\le 6$ & $\le 12$ & $\le 20$ \\
\midrule
\textbf{Moment Gumbel, top 1\,\% trimmed} & \textbf{530} & \textbf{530} & \textbf{530} \\
Moment Gumbel, all maxima & 525 & 530 & 530 \\
Normal distribution & 525 & 530 & 530 \\
Maximum-likelihood Gumbel & 495 & 530 & 530 \\
Maximum-likelihood GEV & -- & 459 & 496 \\
Largest held-out maximum (conformal) & -- & -- & 497 \\
\bottomrule
\end{tabular}
\end{table}

\paragraph{Model or rule?} The simplest rule on the same tokens counts the events of a window whose 4-gram never
occurred in normal behaviour. Scored with the same cross-fit, threshold rule and one-window decision, the model detects
472/530/530/530/531 recordings at budgets of 3/6/12/20/34 normal recordings, the rule none/444/488/497/526. At a budget
of 6 the model detects 86 more recordings (paired bootstrap 95\,\% interval [80, 93]): 53.4\,\% of normal recordings
contain at least one unseen 4-gram, so the rule needs a high threshold, while the model grades \emph{how} unusual a
new gram is. Both are blind to an attacker who uses only seen sequences: removing every event with an unseen gram or
token from each attack recording leaves the model 0/544 (before: 530). This semi-synthetic lower bound says that every
detection on this dataset rests on unseen sequences.

\paragraph{A rate branch.} To target that blind spot we added a branch that scores how much more often a \emph{seen}
gram occurs in a window than in any normal window, $\max_g \log\frac{n_g+1}{m_g+1}$. Alone it detects 225/544 with 7
false alarms (CB-DS 891/891, 3); OR-ed with the model it adds no detection and raises false alarms to 12 (7 with a
0.9/0.1 split of $\alpha'$; CB-DS 4), so it was not joined.

\subsection{RQ3: blind live runs}\label{sec:eval-live}

Two 24-hour blind runs scored all 21 workloads with a campaign of 525 trials each; thresholds were fitted before each
run and fixed during it. The first learnt 12 hours of normal-only capture; the second also learnt the first run's
normal traffic outside its trials, added an operator command to each trial, and sent its alerts through Alertmanager
and the RCA. \Cref{tab:live} gives the primary set.

\begin{table*}[t]
\centering\small
\caption{Token model on the cluster, primary set (6, 12, 24 actions/s), with 95\,\% intervals.}
\label{tab:live}
\begin{tabular}{@{}llll@{}}
\toprule
Metric & Target & Run 1 (12\,h of normal) & Run 2 (more normal) \\
\midrule
Recall & $\ge 95\,\%$ & 315/315 [98.8, 100.0] & 314/314 [98.8, 100.0] \\
Precision (1:1) & $\ge 95\,\%$ & 99.06\,\% [97.26, 99.68] & 100\,\% [98.79, 100.00] \\
Control positives & $\le 5\,\%$ & 3/315 (0.95\,\%) [0.32, 2.76] & 0/314 [0.00, 1.21] \\
p95 kernel $\to$ end of alarm window & $\le 3$\,s & 0.49\,s (p50 0.34\,s) & 0.49\,s (p50 0.35\,s) \\
Harmless binary flagged & reported & 313/315 & 314/314 \\
Operator command flagged & reported & -- & 255/261 \\
Burden (episodes per replica-hour) & low & 0.41 [0.15, 0.73] & 0.17 [0.06, 0.33] \\
Latest run (fresh fit after a cluster change, \cref{sec:eval-v6}) & & \multicolumn{2}{l}{\pending{H0 of pulse-live-v6}} \\
\bottomrule
\end{tabular}
\end{table*}

All four targets are met in both runs. Over all rates, including the slow 1 and 3 actions/s, the runs detect 525/525
and 521/522 attacks (one \texttt{mprotect} trial at 1/s missed) with 5/525 and 0/522 control positives;
\texttt{mprotect} and \texttt{seccomp} appear in the normal behaviour of many workloads and are still caught at every
rate. Scoring adds p50 0.15\,s and p95 0.18\,s per window. Learning more normal traffic lowered the burden from 332
episodes in 817.1 replica-hours to 104 in 630.1, still above a reference budget of 0.1. The two runs are on different
days, so this is not a controlled comparison, but it matches the measured cause: in a sample of the first run, every one
of 146 alarm windows on normal traffic contained an unseen 4-gram, and 89 of 118 episodes came from one Kafka workload.
Held-out 60\,s blocks exceed the threshold 145/48{,}531 times (0.299\,\%, 16.6$\times$ the design level) and, with
more normal data, 350/134{,}905 (0.259\,\%): on the cluster too, $\alpha'$ is a knob. A virtualisation-host stall of
3.5 minutes in the second run produced every cadence issue of that run; excluding the two trial pairs inside it changes
no rate.

\paragraph{The recall is attack-specific.} An alarm could come from the burst of starting a process through
\texttt{kubectl exec} rather than from the attack. Removing from each attack interval every gram that also occurred in
the same trial's harmless-binary and operator intervals leaves 314/314 primary detections (521/522 over all rates), with
the first attack-specific alarm at p50 0.61\,s and p95 0.73\,s after the interval start; no trial is detected only
through launch grams. Counting only alarm windows ending at least 1.5\,s after the start still gives 302/314.

\subsection{RQ4: harmless new processes}\label{sec:eval-benign}

The model flags almost every new process that enters a container through \texttt{kubectl exec}: the harmless binary
in 314/314 primary trials and real operator commands in 419/436 trials over all rates (96.1\,\%). Counting the harmless
binary as a negative, precision on the primary set of the second run would be 50.0\,\%. In a sample of 136 launches of
the harmless binary (134 flagged), alarms came from the launch burst (first alarm after a median 0.44\,s; most
surprising tokens \texttt{mprotect}, \texttt{prlimit64}, \texttt{arch\_prctl}, \texttt{brk}, \texttt{seccomp},
\texttt{keyctl}, \texttt{setgid}, \texttt{setuid}; 61 launches flagged only in the first second) and from the binary's
own loop (1{,}321 of 1{,}323 later alarm windows contain \texttt{clock\_nanosleep}; 73 launches). Neither is malicious.
We tried, each under a registered protocol:
\begin{itemize}
\item \emph{A gate before the model}: an alarm must also be unusual for a pool model of all workloads, or must not fall
  in the 1--2\,s after an \texttt{execve}. Rejected on the datasets: the pool model drops development detections to
  487--509/544; skipping the seconds after \texttt{execve} drops them to 315--365/544 and CB-DS to 549--576/891.
\item \emph{Learning operator activity as normal}: with all other commands learnt, a held-out command is still flagged
  362/439 times (82.5\,\%, against 422/439 with the deployed model), and recall on the primary set falls to 303/314.
  Learning the other folds' operator sessions gives 372/439.
\item \emph{A lineage proxy}: treating the grams the workload model never saw as the session part, the remaining native
  part flags 0/525 idle, 0/439 operator and 0/523 harmless-binary intervals. But a separate session model, thresholded
  by the same Gumbel rule on a few hundred short sessions, sits at 37.7 (42.4 when harmless binaries are learnt too),
  above every measured score, and detects 0/522 attacks.
\end{itemize}
The proxy cannot tell a session gram that the workload also knows from a native one, which is why we built the kernel
lineage tag (\cref{sec:lineage}). Its proof on a test pod passes all eight checks (the container's own loop and the
children of pid~1 stay native, an exec and its children are session, a backgrounded child that outlives the session
stays session), and the token proof still finds one token per call. Runtime exec probes of the kubelet (\texttt{python -c} health checks) are session lineage too, so the
session share of normal traffic differs by workload. \pending{session share per workload from the H1 receipt}
\pending{H1 of pulse-live-v6: native stream on operator and harmless-binary intervals}

\paragraph{A rule after the model.} A rule layer that downgrades an alert whose container only saw runtime execs of
binaries shipped in its image (R1), and labels a non-image binary as drift (R2), removes the operator-command false
alarms of the second run, 256/262 $\to$ 0/262 on the primary set and 422/439 $\to$ 0/439 overall, without touching the
attack recall (314/314), because the campaign's attack binary is not part of any image. This is the benefit side only:
an attack that lives off the land with image tools would also be downgraded.
\pending{H4/H5 of pulse-live-v6: harm of R1 on living-off-the-land attacks; operator tools never seen before}

\subsection{RQ5: cost, alert latency and RCA}\label{sec:eval-cost}

\begin{table}[t]
\centering\small
\caption{Probe overhead per operation, range over three workers (median of 5 rotated rounds).}
\label{tab:overhead}
\begin{tabular}{@{}lrrr@{}}
\toprule
Operation & No probe (ns) & Watched: +ns & Unwatched: +ns \\
\midrule
\texttt{getpid} & 1{,}070--1{,}121 & 407--447 & 212--273 \\
\texttt{read} 4\,KB & 1{,}534--1{,}564 & 453--514 & 214--312 \\
\texttt{futex} & 1{,}155--1{,}181 & 470--535 & 194--259 \\
\texttt{open}+\texttt{close} & 8{,}541--8{,}861 & 1{,}774--2{,}077 & 727--1{,}044 \\
\texttt{mmap}+\texttt{munmap} & 7{,}873--7{,}974 & 1{,}217--1{,}548 & 528--737 \\
\bottomrule
\end{tabular}
\end{table}

\paragraph{Overhead.} \Cref{tab:overhead}: a simple call of a watched container costs 0.4--0.5\,\textmu s more, an
open plus close 1.8--2.1\,\textmu s (the dentry walk). Calls of unwatched containers still pay 0.2--0.3\,\textmu s,
because the tracepoints run for every call of the node before the cgroup filter. The baseline is already high
(\texttt{getpid} about 1.1\,\textmu s), likely from virtualisation and the other probes on the nodes, so absolute
nanoseconds compare better than percentages. The worker unit (loader and Python worker, including a shadow count model
kept for comparison) used 7\,h\,05\,min to 7\,h\,28\,min of CPU per 24\,h, 0.30--0.31 cores of a 24-core node. The
lineage tag adds no cost we can resolve: in 12 cells (three workers, four operations including fork plus exit plus
wait) every 95\,\% interval of the paired ratio against the previous build contains zero except one negative; a fork
with exit and wait takes 332--385\,\textmu s on these VMs. Not measured: the cost per request of a real application.

\paragraph{Gram-map capacity.} With a load generator that makes every call a new gram in one cgroup, the map first loses
grams at 400{,}000 calls/s; at 300{,}000 calls/s (142{,}610 distinct grams per window) it loses none, and the loader
drains it in 0.33\,s (at most 0.41\,s); at 400{,}000 calls/s the longest drain is 0.63\,s, beyond the 0.5\,s frame.
The heaviest node window of the second live run held 6{,}901 distinct grams, about 20$\times$ below the first loss.
When the map is full, grams are dropped and counted by the loader, but the worker does not know the window is
incomplete, which a heavy-load attacker could exploit (\cref{sec:discussion}).

\paragraph{Alert latency.}\label{sec:eval-alert} From the attack's first kernel event to Alertmanager the push bridge
delivered 521/522 alerts with p50 0.51\,s and p95 0.68\,s (per-rate p95 0.65--0.69\,s). A Prometheus scrape path in
parallel needed p50 5.29\,s at a fast interval and p50 60.3\,s at the default one, in time for only 143/522 trials.
Outside the host stall all 2{,}075 alert episodes of the run reached Alertmanager (decision to receipt p50 0.034\,s,
p95 0.064\,s); during the 3.5-minute stall 113 of 190 were lost, because the bridge retries three times without a
durable queue.

\paragraph{RCA.} 467/522 attacks got their own case (89.5\,\%), and the attack binary ranked first in 464/467 cases
(99.4\,\%); a case was ready p50 2.33\,s, p95 2.61\,s after the alert. Of the 55 attacks without their own case, 40
were re-runs of the same binary on the same container merged into an earlier case by design, 14 lost a race (the model
alarmed on the launch burst about 0.2\,s after the start, before the attack binary's exec, and the RCA only considered
execs up to the end of the alarm window), and one had no decision. Rank 1 is helped by the campaign order, which always
ran the attack after the harmless binary. \pending{H6 of pulse-live-v6: random step order and the exec recheck}

\subsection{Model aging and the latest run}\label{sec:eval-v6}

Between the second and the latest run the cluster's nodes were resized and the cluster recovered (Redis, its sentinels
and several services were recreated). Under a registered check, the second run's models and thresholds, unchanged,
scored a new 12-hour normal capture: 278/47{,}082 container blocks exceeded the threshold (0.59\,\%, 95\,\% interval
0.53--0.66\,\%), against 350/134{,}905 (0.26\,\%, 0.23--0.29\,\%) on their own held-out blocks, about 2.3$\times$. The
excess sits in a few workloads (Redis 83/2{,}178, 3.81\,\% against 0.26\,\%; then the API gateway, Kafka and an
inventory service), while others went down. A count model learnt before a cluster change therefore raises the false
alarms of some workloads after it; the latest run uses a fresh fit.
\pending{pulse-live-v6: H0 replication, H3 dilution at PAD 30 and 300, burden}

\subsection{LID-DS 2019}\label{sec:eval-2019}
\pending{lidds2019-confirmation-v2: one run after the latest live run}
